\section{Geometric covers and conditional estimates for the upper bound}
\label{app:upper-details}

This appendix proves the local assertions used in
Proposition~\ref{upper:established-inputs} and the ordinary root estimate
of Proposition~\ref{upper:root-contract}. All clocks and all eligibility
sets in these arguments retain their original block labels.

\subsection{Rectangles and deterministic sums at stopping locations}
\label{updetail:rectangles}
Fix a size-truncated squarefree tuple as in Lemma~\ref{artrunc:lemma}
and a genuine conditional fiber of its label law. For every exterior
name $c$, let $E_c$ be its original deadline. In physical logarithmic
coordinates, a lawful cut satisfies $0\leq\tau_c\leq E_c$.
The class $Q_\tau$ reveals name $c$ at points to the right of $\tau_c$;
all unrevealed eligible names remain in the compound category. Whole
batches at an endpoint are treated on the prescribed side. Let $\pi$
be the original categorical product law on this fiber and set
\[
 B=\sum_cE_c,\qquad
 \Xi_\tau=\sum_c\tau_c-B-\log\pi(Q_\tau),\qquad
 Z(t)=1+\sum_{x\leq t}e^{x-t}.
\]
The sum uses all original prefix occurrences, including those that may
not enter a selected descendant list. Revealed factors translate the
remaining boxes. The unresolved displacement in coordinate $c$ is at
most $C_He^{\tau_c}Z(\tau_c)$. Bounding the remaining child coordinates
by their original size rectangles therefore gives
\begin{equation}
 \operatorname{Vol}_{\rm norm}(Q_\tau)
 \leq C_H\pi(Q_\tau)e^{\Xi_\tau}\prod_c Z(\tau_c).
 \label{updetail:rectangle}
\end{equation}
In an ancestor atom $\alpha$, this formula retains its absolute
probability $\pi(\alpha)$ and its inherited geometric multiplier.
Conditioning divides by $\pi(\alpha)$ only temporarily; reinsertion
restores it.

The cuts used for moment estimates are fixed before drawing the
occurrence process. To handle first failures, partition a deterministic
physical interval $[a,b]$ into cells of length at most $\Delta$, with a
bounded number of endpoints per unit interval. If $v$ is the right
endpoint of a cell and $t$ belongs to that cell, positivity gives
\begin{equation}
 Z(t)\leq e^\Delta Z(v).
 \label{updetail:mesh-shot}
\end{equation}
Indeed every term present at $t$ is also present at $v$, and its kernel
changes by $e^{v-t}$. Suppose first-failure atoms form a lawful antichain
and each has capacity at most
\[
 C\pi(q)e^{-h}\omega(t)Z(t)^d\prod_c Z(\tau_c),
 \qquad d\leq H,
\]
where $\omega(t)=(1+\operatorname{dist}(t,\{a,b\}))^{-A}$, or its
one-ended version, and the outside cuts $\tau_c$ are deterministic.
Within one mesh cell $\omega(t)\leq C_{\Delta,A}\omega(v)$ and the
sum of the atom probabilities is at most one. Thus the total failure
capacity is bounded by $Ce^{-h}S$, with
\[
 S=\sum_v\omega(v)Z(v)^d\prod_c Z(\tau_c).
\]
There is no occupancy multiplier in this sum: the stopped atom
probabilities have already been summed.

For a Poisson process with bounded intensity per unit interval,
$Z(v)$ has all fixed moments uniformly in deterministic $v$. This
follows either by its Laplace functional or by splitting into unit
intervals and using the exponentially decreasing coefficients.
For sufficiently large fixed $A$, Minkowski's inequality and
$\sum_v\omega(v)<\infty$ give $\mathbb E S^p\leq C_{H,p}$ for
each required fixed $p$. Finite sums over the endpoint contexts have
the same property. A cut within bounded distance of an endpoint is
covered by finitely many extra mesh cells. An incoming selected-list
sum is bounded pointwise by the original-list sum at the same
deterministic threshold.

The probability tests themselves may be made on the complete
deterministic mesh, including empty cells. The risk has one-sided
jumps and bounded linear speed; survival on that mesh implies the
corresponding continuous necessary wall after a fixed reserve
increase, using the appropriate side for the jump convention. The
rising wall has bounded derivative on the fixed interval, so the
same conclusion holds for its first-hit strip. Consequently the
fixed-time bank estimates below are only applied at deterministic
mesh locations.

Under insertion of $q$ fixed original occurrences, every deterministic
$Z(v)$ increases by at most $q$. Keeping the mesh and contexts fixed
therefore gives
\begin{equation}
 S(G+\text{insertions})\leq C_{H,q}S(G).
 \label{updetail:shot-insertions}
\end{equation}
The categorical log-likelihood changes at each fixed test by at most
$C_Hq$, which is paid by the reserve. A Poisson--Mecke expansion of a
fixed shot monomial uses finitely many partitions of its occurrences;
apply this insertion bound in each term before summing the deterministic
locations. This proves the fixed-moment and insertion assertions needed
for all the actual stopping covers. In particular it does not estimate
a shot directly at an unrestricted data-dependent location.

\subsection{Positive prime expansion with the original root}
\label{updetail:prime-expansion}
Let $\mathcal P_i$ be the original prime ranges, including $2,3$ in
$\mathcal P_1$, and put
\[
 H_i=\sum_{p\in\mathcal P_i}\frac1p,\quad
 B_{\rm geo}=\sum_iV_i=\sum_i(h_i-1)L_i.
\]
Take fixed positive $Z_i,s_i$, put $s=s_1>0$, and define
\[
 J=\sum_iL_i+\sum_i s_iV_i-\sum_iZ_iL_i,
 \qquad
 X(G)=\frac{\sum_iN_i(G)\log Z_i-\sum_i s_iV_i}{s}.
\]
Let $\mathbb P_Z$ give independent occurrence counts
$\operatorname{Pois}(Z_i/p)$ at each $p\in\mathcal P_i$.
For every nonnegative symmetric functional $A$ on occurrence lists,
Tonelli and the Poisson exponential formula give the exact identity
\begin{align}
 &e^{-\sum_iL_i}
 \sum_{\mathbf n\geq0}\prod_i\frac1{n_i!}
 \sum_{p_{i,a}\in\mathcal P_i}
 A(G)\prod_{i,a}\frac1{p_{i,a}}
 \notag\\[-2pt]
 &\hspace{25mm}=
 e^{-J+\sum_iZ_i(H_i-L_i)}
       \mathbb E_Z[e^{-sX(G)}A(G)].
 \label{updetail:exact-expansion}
\end{align}
To verify the exponent, the unnormalized list sum equals
$e^{\sum_iZ_iH_i}\mathbb E_Z[A\prod_i Z_i^{-N_i}]$, and
$\prod_i Z_i^{-N_i}=e^{-sX-\sum_i s_iV_i}$.

For the saddle, $\log Z_i=G_i$ and $X$ is exactly the uncolored
routing root. At node $j$, multiplying minus the local log-probability
by $s_j$ gives $G_j$ when the occurrence exits and $G_j-G_{j+1}$ when
it enters the child. Along each route these contributions telescope
to $G_i$ at its original block. Subtracting the physical budgets
proves the identity without a stationarity assumption or resampling
any child list.

Define $I_0(G)=e^{-B_{\rm geo}}L^{(k+1)}(\mathbf a(G))$ on the
squarefree size-truncated support, and $I_0=0$ elsewhere. Each squarefree
selection occurs $n_i!$ times among ordered lists, so
Lemma~\ref{artrunc:lemma} and \eqref{updetail:exact-expansion} imply
\begin{equation}
 K_k(\ell)\leq
 2e^{-J+\sum_iZ_i(H_i-L_i)}
       \mathbb E_Z[e^{-sX}I_0(G)].
 \label{updetail:arithmetic-upper}
\end{equation}
The geometric caps hold first on that actual support. If they give
$I_0\leq C C_l(G)$ there, with each $C_l$ nonnegative on every list,
then $I_0\leq C\min_l C_l$ on all lists, because of the zero
extension. Thus unrestricted repeated lists require no artificial
size-truncation assertion. Their numerical cap minimum is evaluated
under the same atomic Poisson law and the same root.

The uniform prime harmonic estimate gives $|H_i-L_i|\leq C_H$;
the discrepancy in \eqref{updetail:arithmetic-upper} is bounded since
$Z_i\leq H$. Deleting a fixed initial prime set pays its finite Euler
product and changes each $H_i$ by only a fixed amount. For later cell
completion, if $\mu_c^*\geq\mu_c$ and
$\sum_c(\mu_c^*-\mu_c)\leq C_H$, the product Poisson density ratio is
at most $e^{C_H}$. It is applied once to the whole nonnegative
integrand, retaining its minimum, powers and root. A change of the
physical cell budget is compensated in the external root translate.

\subsection{Separation of the rank intervals of rare comparisons}
\label{updetail:rank-separation}
In one maximal owner, write $\varepsilon=1-\sigma>0$ and
\[
 \kappa_r=(r+z)^{-\varepsilon},\qquad
 g_r(a)=a-\kappa_rE_a(z).
\]
The transported old coefficients are retained individually. They are
no larger than the natural coefficient at the first native capacity.
Every lawful complete reference history has nonnegative cost by the
KKT and reference-compression results of
Appendix~\ref{app:endpoint-details}.

Choose a fixed sufficiently small $\epsilon_0(H)>0$. An edge $r\to q$
is weak when $|g_r(r)-g_r(q)|\leq\epsilon_0$. If $m_-,m_+$ are the
two minimum reference costs, their common increment is
$m_+-m_-=L_i\nu_i$. Both are nonnegative and the endpoint scores
dominate $1+m_-$ and $1+m_+$. Hence
\[
 |\nu_i|L_i\leq m_-+m_+\leq X_i+Y_i.
\]
Since the natural mean is comparable to $L_i$, sufficiently small
$(1+X_i)(1+Y_i)/(1+\mu_i)$ implies this fixed weakness condition.

We record the three strict comparisons needed below. The normalized
secant estimate of Appendix~\ref{app:endpoint-details} is
\[
 \beta_b(y)=\frac{\log((y\log y-b\log b)/(y-b))}{\log y},
 \qquad \beta_b'(y)\leq-\frac1{24H^2}.
\]
For a weak edge $r\to q$ and an earlier capacity $h>r$, put
$T_a=(a+z)^{-\varepsilon}(E_a-E_q)/(a-q)$. Weakness gives
$|T_r-1|\leq\epsilon_0/(r-q)$, whence
$\beta_{q+z}(r+z)-\varepsilon\leq2\epsilon_0/\log2$.
If $\epsilon_0\leq\log2/(192H^2)$, then
$\log T_h\leq-\log2/(48H^2)$. At the other endpoint,
\[
 g_h(r)-g_h(q)
 =(r-q)\{1-(\kappa_h/\kappa_r)T_r\},
 \quad \kappa_h/\kappa_r\leq e^{-\varepsilon_{\min}/H}.
\]
Taking also
$\epsilon_0\leq(e^{\varepsilon_{\min}/H}-1)/2$ gives a positive
lower bound at both $r$ and $h$. Concavity then proves
\begin{equation}
 g_h(a)-g_h(q)\geq\delta_H>0\qquad(r\leq a\leq h).
 \label{updetail:earlier-rank}
\end{equation}
Reducing the coefficient increases this difference, so it also holds
in the relevant old groups. If $q\geq1$ and $0\leq a<q$, strict
concavity $-g_r''\geq H^{-2}$ gives
\[
 g_r(r)-g_r(a)\geq-H\epsilon_0+
                \frac{(q-a)(r-a)}{2H^2}
 \geq\frac3{4H^2}
\]
after $\epsilon_0\leq1/(4H^3)$. Finally the forward secant comparison
gives, for $q<h<r$,
\begin{equation}
 g_h(q)-g_h(h)\geq\gamma_H>0.
 \label{updetail:intervening-rank}
\end{equation}
All thresholds are chosen before any lengths.

Suppose two weak edges have $r_1>r_2$ but $q_1<r_2$. Output ordering
gives $q_2<q_1$, so $q_1\geq1$. In any left completion for the second
edge, all earlier retained ranks are at least $r_2$. Lower every such
rank to $q_2$, keeping the whole later history. This is a legal
complete history. Nonnegativity and \eqref{updetail:earlier-rank}
therefore give $X_2\geq1+\delta_HL_1$.

For a right completion of the first edge, fix the prefix through its
block and the suffix after the second block. Between their right
endpoints $b_1,b_2$, all capacities exceed $q_1$, so none of the
remaining names expires there. The reference cost, as a function of
their ordered cut times, is a constant plus
\[
 \sum_{a=k_{\rm end}+1}^{q_1}
 \int_{b_1}^{\tau_a}\{g_t(a)-g_t(a-1)\}\,dt.
\]
Each summand is concave because $\kappa(t)$ is nondecreasing. Its
minimum on the ordered simplex occurs at a vertex, with all cuts
at $b_1$ or $b_2$. Thus one constant rank $a\leq q_1$ can be used
on that interval without increasing the cost. If $a<q_1$, lower the
rank on the first block to $a$ and use the strict concavity comparison
above: its original cost was at least $c_HL_1$. If $a=q_1$, replace
the intervening ranks by their full eligible capacities and use
\eqref{updetail:intervening-rank}: the original cost was at least
$c_HL_2$. Both comparisons preserve the complete outside history.
Consequently
\[
 Y_1\geq1+c_H\min(L_1,L_2),\quad
 \overline T_1\geq c_1\frac{1+\min(L_1,L_2)}{1+L_1},\quad
 \overline T_2\geq c_2\frac{1+L_1}{1+L_2}.
\]
If $L_2\geq L_1$ the first ratio is bounded below; otherwise their
product is. Hence two sufficiently rare edges must have $q_1\geq r_2$.
Applying this to every pair proves disjointness of the rank intervals
$(q_i,r_i]$. The argument also applies to terminal nonzero edges:
$z=1$, but their secant bases are at least two and their marks remain
nondegenerate. The terminal zero-output edge remains ordinary.

\subsection{The two scales of a proper zero-output comparison}
\label{updetail:zero-scales}
Consider a proper zero-output edge with sufficiently small persistence ratio $\overline T$. Let $M$ be the selected native block, $A,B$ the preceding and following
native lengths, and $V_{\rm old}$ the complete old length. Put
\[
 P=V_{\rm old}+A,\qquad T=P+M+B,\qquad
 R_{\rm exit}=\max\{R_0,\min(\sqrt{1+A+M+B},d^{-1})\}.
\]
Full input and zero output force the two completion times to be
$P$ and $P+M$ in every outside context. The clipped-prefix estimate
from Appendix~\ref{app:endpoint-details} gives
$\overline C_v\geq c_H\min(t_v,T-t_v)$.
If the block fails to cross $T/2$ strictly, one of its endpoint costs
is at least $c_HM$ and the persistence ratio is bounded below.
Thus sufficiently small ratio forces $P<T/2<P+M$.

For upper bounds on the minima, retain every earlier name to its own
deadline and remove the current names at the selected endpoint. Both
contexts have no good groups. Complete owner balance gives
\[
 \overline C^-_{\rm full}
 =\sum_{\rm old}L_gg_g(r_{\rm top})+
   \sum_{\rm native\ before}L_gg_g(r_g),\qquad
 \overline C^+_{\rm full}
 =-\sum_{\rm native\ after}L_gg_g(r_g).
\]
Their costs are nonnegative and at most $C_HP$ and $C_HB$,
respectively. Together with the clipped lower bound and the early/late
score definition, this proves
\begin{equation}
 X\asymp_H1+V_{\rm old}+A+R_{\rm exit},\qquad
 Y\asymp_H1+B.
 \label{updetail:intrinsic-zero-pair}
\end{equation}
In particular rarity implies $V_{\rm old}+A+B\leq C_H\tau M$;
for a small fixed $\tau$, $M$ dominates the entire old and native
length. Uniformly even as $d\downarrow0$,
\[
 R_{\rm exit}\asymp_H
 \varrho:=\max\left\{\varrho_0,
 \frac{c_0\sqrt{1+M}}{1+d\sqrt{1+M}}\right\}.
\]
Conversely a sufficiently small
$(1+B)(1+V_{\rm old}+A+\varrho)/(1+M)$ forces that dominance and
midpoint crossing, so the scales characterize the same rare region.
For the original mark
$W=\log(r+z)-\log z\,1_{\rm child}$, complete balance also gives
\[
 eM=\sum_{\rm old}L_gg_g(r_{\rm top})+
     \sum_{\rm native\ne i}L_gg_g(r_g),
 \qquad e=(r+z)^\sigma\mathbb EW-r,
\]
and therefore $|e|M\leq C_H(V_{\rm old}+A+B)$.
The two reserves in \eqref{updetail:intrinsic-zero-pair} remain
separate when applying a one-end or two-end probability estimate.

\subsection{Necessary one-end tests on original native intervals}
\label{updetail:one-end-banks}
Consider a weak edge $r\to q$ with $q\geq1$ on $[a,b]$, of length
$L$, in the rare region $\overline T<\tau$, where $\tau>0$ is
sufficiently small in terms of $H$. Write $S=r+z$, $B=q+z$, $\nu=g_r(r)-g_r(q)$, and let
$\mathcal J$ consist of the later native groups with $q<h<r$.
Their total length is $M_{\mathcal J}$.
We first prove
\begin{equation}
 1+X\geq c_H\sqrt{1+L},\quad
 1+Y\leq C_H\tau\sqrt{1+L},\quad M_{\mathcal J}\leq C_HY.
 \label{updetail:one-end-scales}
\end{equation}
At a left completion the current block retains a proper deepest
subset. If the completion is late, this is a good group and contributes
$\sqrt L$; if it is early, its last positive time is at least $b$,
and clipped-prefix charging gives a cost at least $c_HL$.
This proves the first inequality, and rarity gives the second.
For any right completion, compress the cuts between $b$ and the end
of $\mathcal J$ to the two endpoints, by the same concavity argument
as in Section~\ref{updetail:rank-separation}. The compressed rank is
some $a'\leq q$. If $a'<q$, lowering the current block costs at least
$\delta_HL$; if $a'=q$, filling the intermediate groups costs at least
$\gamma_HM_{\mathcal J}$. Complete reference nonnegativity thus gives
$Y\geq\min(\delta_HL,\gamma_HM_{\mathcal J})$. For sufficiently small
$\tau$, rarity excludes the first alternative, proving the last claim.
These arguments also hold for $z=1$ and positive $q$.

The endpoint reference minima differ by $\nu L$, so
$|\nu|L\leq X+Y$. Choose the deterministic interval length
\begin{equation}
 \ell=c_{\rm bank}\frac{L^2}{(1+X)^2}.
 \label{updetail:effective-length}
\end{equation}
By choosing $c_{\rm bank}$ sufficiently small and then $\tau$
sufficiently small, \eqref{updetail:one-end-scales} and rarity imply
\[
 \ell_{\rm scalar}\leq\ell\leq L/16,
 \qquad |\nu|\sqrt\ell\leq C_H,
 \qquad
 \frac{1+Y}{\sqrt\ell}\asymp_H
 \overline T:=\frac{(1+X)(1+Y)}{1+S^\sigma L}.
\]
Indeed rarity gives $1+X\leq C_H\tau L/(1+Y)$ and hence
$\ell\geq c_H\tau^{-2}$. Neither $\ell$ nor its location is selected
from the sampled word.

Group the deepest $q$ names and the true compound child into the
category $B$, keeping the other names separate. Reveal the complete
coarse word in $\mathcal J$, and put
\[
 H_{\mathcal J}=\sum_{q<h<r}
 \{N_h\log(h+z)-K_{B,h}\log B-(h-q)L_h\}.
\]
Read the current interval backwards and set
\[
 W=\log S-\log B\,1_B,\qquad
 C(t)=H_{\mathcal J}+\sum_{x<t}W_x-(r-q)t.
\]
The original rate is $S^\sigma$ and its mean drift is $-\nu$.
Earlier exterior coordinates keep their full original sizes;
the intermediate ones keep their shorter deadlines; the $B$
coordinates are projected into their original full rectangle.
The unresolved exterior budget is the full budget minus
$(r-q)t+\sum_{\mathcal J}(h-q)L_h$, while the observed negative
log-probability is the mark sum above. Formula
\eqref{updetail:rectangle} therefore gives the actual partial capacity
$C_H\pi_t e^{C(t)}\operatorname{Shot}_t$.
Stop at the first mesh failure of
$C(t)\geq-h-A\log(1+t)$, including $t=0$.
Section~\ref{updetail:rectangles} bounds the union of failures by
$CS e^{-h}$ with all required moments.

Because $H_{\mathcal J}\leq(\log H)N_{\mathcal J}$ pointwise,
every survivor satisfies the necessary test
\begin{equation}
 \sum_{x<t}W_x-(r-q)t
 \geq-h-C_HN_{\mathcal J}-A\log(1+t),\qquad 0\leq t\leq\ell.
 \label{updetail:necessary-one-end}
\end{equation}
First make this implication with the entire intermediate word fixed,
then release that word with its original mass. The right side now
reads only $N_{\mathcal J}$; any future private test remains in the
joint surviving-word operator.

Let $Q_h$ be the conditional private probability of
\eqref{updetail:necessary-one-end}. The signed logarithmic-wall
one-end estimate, on the original clock with bounded scaled drift,
gives
\[
 \mathbb E_{G_i}Q_h^\sigma
 \leq C_H\min\left\{1,
 \left(\frac{1+h+C_HN_{\mathcal J}}{\sqrt\ell}\right)^{2\chi(\sigma,\rho)}
 \right\},\qquad
 \rho=\frac{(\mathbb EW)^2}{\mathbb E W^2}.
\]
The intermediate count is independent of this original native clock,
with mean at most $C_HM_{\mathcal J}\leq C_HY$. Its fixed moments
are at most $C_p(1+Y)^p$. Integrating over all its counts proves
\begin{equation}
 \mathbb E Q_h^\sigma\leq C_H(1+h)^K\overline T^{2\chi(\sigma,\rho)}.
 \label{updetail:one-end-moment}
\end{equation}
To compare with the exact-type parameter, write $a_0=\log S$,
$b_0=\log B$, $c=r-q$, $\lambda=S^\sigma$, $\mu=\lambda L$,
and $p=B/S$. Its centered type parameter and correlation are
\[
 \bar p=(a_0-c/\lambda)/b_0,
 \qquad
 \bar\rho=\frac{(cL/(\mu+1))^2}
 {(cL/(\mu+1))^2+b_0^2\bar p(1-\bar p)}.
\]
The identity $\bar p-p=-\nu/(b_0\lambda)$ gives
$|\bar\rho-\rho|\leq C_H(|\nu|+1/L)$.
As $\overline T\geq c_H/\sqrt\ell$,
$|\log\overline T|\leq C_H+\tfrac12\log\ell$.
The scaled-drift bound and $\ell\leq L/16$ bound the product of
this logarithm and the correlation error. Compact local Lipschitz
continuity of $\chi$ therefore permits either parameter in
\eqref{updetail:one-end-moment} at fixed comparison cost.
This includes every count and all $n+1$ conditional spacings.
Fixed insertions change the reserve by at most a constant per point;
the Mecke argument of Section~\ref{updetail:rectangles} attaches
the required shot powers at the same exponent $\sigma$.

For a proper $r\to0$ edge, assume explicitly
$1+X\geq c_X\sqrt{1+L}$. The child weight satisfies
$z\geq2^{u_{\min}}>1$, so the two marks remain nondegenerate.
Here $\mathcal J$ consists of all later native groups. At a right
completion, postpone the current cuts to their individual deadlines.
The complete earlier history is unchanged, and the difference of
reference costs is
$\sum_{0<h<r}[-g_h(h)]L_h\geq\gamma_HM_{\mathcal J}$.
Thus $M_{\mathcal J}\leq C_HY$ directly. The preceding proof applies
with $B=z$, $c=r$, and the child category in place of the deepest
suffix. In particular it supplies precisely the same moment
\eqref{updetail:one-end-moment}. The left-score hypothesis is needed:
there is no current proper-good group when the output is zero.

The sets consisting of a current block and its intermediate native
groups are disjoint for distinct rare edges of an owner, by rank
separation, and are disjoint across owners. Their label intervals
have length at most $L/16$ and avoid a count middle in the same
block. If the count middle lies in an intermediate group, its weight
is decreasing in that count and the necessary reserve is increasing.
For independent copies $N,N'$,
\[
 2\operatorname{Cov}(f(N),g(N))
 =\mathbb E[(f(N)-f(N'))(g(N)-g(N'))]\leq0
\]
when $f$ is decreasing and $g$ increasing. Applying this with all
other groups fixed supplies the required separation without an
independence assertion for the shared count.

\subsection{Leaving a terminal suffix unread}
\label{updetail:terminal-trimming}
All terminal binary weak blocks lie in the negative ordinary-drift
head. The ordinary-suffix geometry gives, before trimming,
\[
 D_0\geq b_HV_{\rm old}+g_0L_{\rm head},\qquad
 U\geq c_HT_{\rm total},\qquad M_*\geq c_HT_{\rm total}/H.
\]
A terminal necessary test may read intermediate counts within this
suffix. If so, start the suffix at the first capacity at most its
output $q$. The removed initial segment has length
\[
 \delta\leq M_{\mathcal J}\leq C_HY
 \leq C_H\tau\sqrt{1+L}\leq C_H'\tau L.
\]
For several tests use the latest one; rank separation puts all earlier
read sets before it. The old head contains its current length $L$.
Choosing the fixed $\tau$ small enough, bounded ordinary drifts give
$E_H\delta\leq D_0/4$, $\delta\leq D_0/4$ and $\delta<M_*/2$.
The largest original root block is therefore retained and
\[
 U'=U-\delta\asymp U,\quad M_*'=M_*,\quad
 A_*'=A_*-\delta,\quad B_*'=B_*,\quad
 D'=D_0-\sum_{\rm removed}e_hL_h\asymp D_0.
\]
The new suffix still has $D'\geq c_HV_{\rm pre}'$ and
$U'\geq c_HT_{\rm total}$. Moreover $1+D'+A_*'$ is comparable to
$1+D_0+A_*$, so its ordinary factor is unchanged up to constants.
Every clock read by the necessary masks now belongs to the actual
prefix of this suffix. When no terminal binary rare edge is present,
no trimming is needed.

\subsection{Whole-band estimates with two variable reserves}
\label{updetail:whole-band}
Fix a compact nondegenerate positive binary law $W$, a compact positive
rate $\nu$ and speed $c$, and a power $\sigma$ in a compact subinterval
of $(0,1)$. Suppose $M,L,R\geq1$ and the band width $r$ satisfies
\begin{equation}
 1\leq r\leq C\sqrt M,\qquad r\leq CR,\qquad
 |\nu\mathbb EW-c|M\leq C(L+R).
 \label{updetail:band-assumptions}
\end{equation}
All interval locations are deterministic. Use a middle interval of
length comparable to $M$, disjoint from the retained end intervals.
An additional later-side count $N_B$ is independent, with Poisson mean
at most $CL$. Let the positive left test have reserve at most
$C(1+h+N_B)$ and the reflected negative right test have reserve at most
$C(1+h+R+(1+|j|)r)$. Both allow a fixed logarithmically lowered wall.
After fixing the entire outside private word, assume the whole band
permits at most $Cr$ values of the middle binary count.

Put $T=LR/M$ and
$\rho_{\rm nat}=(\mathbb EW)^2/\mathbb EW^2$.
There is a numerical majorant $B_{h,j}(G)$ for the conditional
probability of the band and its retained tests such that
\begin{equation}
 \mathbb E B_{h,j}^{\sigma}
 \leq C(1+h)^K(1+|j|)^K
       (r/\sqrt M)^\sigma
       \min\{1,T^{2\chi(\sigma,\rho_{\rm nat})}\}.
 \label{updetail:whole-band-moment}
\end{equation}
Here the entire permitted count band is summed before the power.
To prove this, if $\max(L,R)\leq C_0\sqrt M$, retain end intervals
of length $M/16$. Otherwise retain only the smaller-reserve end,
of length $\ell=\min(M/16,c_*M^2/\max(L,R)^2)$.
Conditioning on all outside labels, the middle count is binomial,
so its maximum atom gives the pointwise majorant
\[
 B_{h,j}=\min\{1,Cr/\sqrt{1+N_{\rm mid}}\}\,q_{\rm left}q_{\rm right},
\]
with an omitted end factor equal to one. The bound is uniform in the
outside-dependent center. Release those outside labels only afterwards.
The Poisson negative moment supplies $C(r/\sqrt M)^\sigma$, including
small middle counts. On every retained end the scaled drift is bounded
by \eqref{updetail:band-assumptions}. Thus the signed one-end estimate gives
\[
 \mathbb E q_H^\sigma\leq C
 \min\{1,[(1+H)/\sqrt{1+\ell}]^{2\chi_c}\},\quad
 \chi_c=\chi\left(\sigma,
 \frac{(c/\nu)^2}{(c/\nu)^2+\operatorname{Var}W}\right).
\]
If both ends remain, independence and the moments of $N_B$ give the
path factor $\min(1,(LR/M)^{2\chi_c})$ times the displayed reserve
polynomials. The condition $r\leq CR$ pays all translated right bands.
If only one end remains, its reserve divided by $\sqrt\ell$ is bounded
by those polynomials times $LR/M$, yielding the same result.
If $\ell$ is bounded, $\max(L,R)\geq cM$ and $T\geq c$; discard
the path test at fixed cost while retaining the middle atom.
On $T<1$ the drift hypothesis implies
$|\chi_c-\chi(\sigma,\rho_{\rm nat})|\leq C(L+R)/M\leq2CT$.
The replacement costs at most $e^{C|T\log T|}$. This proves
\eqref{updetail:whole-band-moment}; for $T\geq1$ both numerical path
factors equal one.

For actual prime clocks these assertions follow by a pointwise
quantile comparison. On a deterministic physical interval of length
$\ell$, its harmonic cumulative mass is $t+O_H(1)$ for every initial
length $t$, with the same statement for half-open intervals and their
reflections. The harmonic-to-physical quantile map therefore satisfies
$|T(s)-s|\leq C_H$. Its image of a homogeneous Poisson process is
the required atomic prime process, with repeated occurrences allowed.
A positive-jump wall is tested at its pre-jump minimum; a negative
wall at its post-jump minimum. During the collapse of any atom the
homogeneous partial sums dominate these minima up to the bounded
clock error. An actual necessary wall is consequently contained in
the homogeneous wall with a fixed reserve enlargement, for every mark
word and every atom multiplicity. Logarithmic allowances change by
only a fixed amount.

For clarity, the full-count homogeneous estimate used here follows
from Theorem~\ref{pr:one}. Let $m=\lfloor\nu\ell/2\rfloor$.
The probability of fewer than $m$ arrivals by $\ell$ is exponentially
small. Otherwise survival implies a rank wall through the first $m$
independent exponential gaps. For the positive sign, survival at the
$j$th arrival implies
$cT_j\leq w_+j+h+A\log(1+T_j)$ and hence
$T_j\leq C(j+h+1)$. For the negative sign, when this latter bound
fails the signed risk is already nonnegative; in the complementary
case the same logarithmic comparison applies. In either direction
the rank-wall reserve is $C(1+h)$ with a fixed logarithmic allowance.
Theorem~\ref{pr:one} gives
$C(1+h)^{2\chi}(1+\ell)^{-\chi}$, and the exponentially small
exception is absorbed. This proof includes all original counts.

The fixed-width version is especially useful at first hits. Split
a deterministic interval $[0,z)$ into its left quarter, middle half,
and right quarter. Given both end words, a fixed terminal band admits
at most a fixed number of middle low counts. Hence, uniformly over
every band center $b$,
\[
 \sup_b q_b(G)\leq C(1+N_M)^{-1/2}\ell_h(G_L)r_{r_0}(G_R),
\]
and independence of these three deterministic restrictions gives
\begin{equation}
 \mathbb E(\sup_bq_b)^\sigma
 \leq C(1+h)^{2\chi}(1+r_0)^{2\chi}
                   (1+z)^{-\sigma/2-2\chi}.
 \label{updetail:fixed-band}
\end{equation}
There is no extra terminal-height factor. The estimate is uniform in
$z$ in a prescribed comparable range, not an estimate for a supremum
over $z$; stopping covers retain their deterministic summation weights.

Insertion of $q$ fixed occurrences decreases the middle negative
count factor, relaxes the positive left reserve by at most $qw_+$,
and makes the negative right wall stronger. The same middle-binomial
argument can leave the original middle marks free after conditioning
on the inserted marks. Thus both displayed bounds remain valid with
that left reserve shift. To attach shots, for
$\beta=4/\sigma_{\min}$ form
\[
 C_z(G)=\sum_{h,r_0\geq0}
 (1+h)^{-\beta}(1+r_0)^{-\beta}
 (1+N_M)^{-1/2}\ell_h(G_L)r_{r_0}(G_R).
\]
Fractional subadditivity and \eqref{updetail:fixed-band} show
$\mathbb EC_z^\sigma\leq C(1+z)^{-\sigma/2-2\chi}$, since each
reserve exponent is at most $-3$. Reindexing after insertion gives
$C_z(G+q\text{ points})\leq C_q C_z(G)$.
Expanding an integrably weighted fixed-degree shot polynomial into
Poisson occurrence partitions and applying Mecke now preserves this
power of $1+z$. The expansion is valid for atomic intensities as well.

For the proper zero-output application, take the separate scales
$L=1+B$ and $R=1+V_{\rm old}+A+\varrho$ from
\eqref{updetail:intrinsic-zero-pair} and band width $\varrho$.
They satisfy all of \eqref{updetail:band-assumptions}, including
$\varrho\leq CR$. If $d\sqrt M<1$, then
$\varrho\asymp\sqrt M$ and $R\geq c\sqrt M$; the left-only
calculation gives the same bound with a width factor comparable to
one. No series whose ratio tends to one as $d\downarrow0$ is used.

\subsection{Joint count weights and the conditional routing tower}
\label{updetail:count-tower}
For each proper owner choose a fixed middle in a largest native block,
write $n_j$ for its count, and set
$M_j=(1+d_j\sqrt{1+n_j})^{-1}$.
For each simple edge let $Q_i(h)$ be its necessary-test probability.
The only possible overlap between a count middle and these tests is
when the middle lies in an intermediate set $\mathcal J_i$.
With all other clocks fixed, $M_j^{\sigma_j}$ decreases in $n_j$
and $Q_i^{\sigma_{o(i)}}$ increases. The covariance inequality proved
above extracts $\mathbb EM_j^{\sigma_j}$. Each middle lies in at
most one such set by rank separation. The remaining test inputs are
disjoint original groups, so
\[
 \mathbb E\left[\prod_{j<m}M_j^{\sigma_j}
                    \prod_iQ_i(h)^{\sigma_{o(i)}}\right]
 \leq\prod_{j<m}\mathbb EM_j^{\sigma_j}
       \prod_i\mathbb EQ_i(h)^{\sigma_{o(i)}}.
\]
Splitting a Poisson count below and above half its mean gives
\[
 \mathbb EM_j^{\sigma_j}
 \leq C_H(1+d_j\sqrt{1+\mu_{C_j}})^{-\sigma_j},
\]
uniformly even as $d_j\downarrow0$. Together with
\eqref{updetail:one-end-moment} this gives all proper count and simple
native factors at their respective owner powers. Fixed insertions
only decrease $M_j$ and shift each reserve by a fixed amount; the
same Mecke argument attaches the permitted shots. One unweighted
count insertion pays its original mean.

Here is the finite conditional argument that uses these environmental
costs. Fix the complete original uncolored word $G$ and let
$\mathcal Q_j$ reveal its first $j$ coarse owner routes. In a native
block of owner $j$, every earlier owner has already expired; its
earlier routes are deterministically child. Thus on every genuine
$\mathcal Q_{j-1}$ atom its current labels still have probabilities
$1/(r+z_j)$ for each exterior name and $z_j/(r+z_j)$ for the child,
independently over original slots. Let $E_j$ avoid the count middle
and satisfy
$\pi_G(E_j\mid\mathcal Q_{j-1})\leq p_j(G)$ uniformly on those
atoms; for the terminal owner assume the corresponding conditional
bound on $E_m$. Native-only masks have this property.

For any uncolored-measurable thresholds $\xi_j$, put
$\zeta_j=\eta_j-\xi_j$. Condition on the entire current coarse route
outside its middle. Then $E_j$ is fixed and
\[
 \zeta_j=H_0-\log z_j\,K_j,\qquad
 K_j\sim\operatorname{Bin}(n_j,z_j/(r+z_j)),
\]
where $H_0$ may depend arbitrarily on the actual incoming list.
The binomial maximum atom and the two exponential half-lattice sums
give, uniformly in $H_0$ and also for $n_j=0$,
\begin{align*}
 \mathbb E[e^{\zeta_j}1_{\{\zeta_j\leq0\}}1_{E_j}
                  \mid\mathcal Q_{j-1}]&\leq C M_jp_j,\\
 \mathbb E[e^{-d_j\zeta_j}1_{\{\zeta_j>0\}}1_{E_j}
                  \mid\mathcal Q_{j-1}]&\leq C M_jp_j.
\end{align*}
For the negative side use the geometric sum with ratio $e^{-\log z_j}$;
for the positive side use $e^{-d_j\log z_j}$ and cap the bound by one.
At $d_j=0$ the probability bound suffices. The arbitrary center is
removed before releasing the outside route. Future child words are
not fixed in this calculation.

Backward conditioning now bounds every product of positive half-line
factors followed by one negative half-line factor by
$C^j\prod_{l\leq j}M_lp_l$, and the fully surviving branch by
$C^{m-1}p_m\prod_{l<m}M_lp_l$. For completeness set
$a_j=M_jp_j$ for $j<m$, $a_m=p_m$, $r_j=\sigma_j/\sigma_1$, and
\[
 B_0=\sum_{j\leq m}r_j\log a_j,\quad
 \theta_j=r_j^{-1}\sum_{l>j}r_l\log a_l,\quad
 \xi_j=X/r_j+\theta_j.
\]
Stop the actual projected route at the first $\eta_j\leq\xi_j$.
The original capacity is $C_H\pi_G(q)e^{P_j}$, and the carry identity
gives on this branch
\[
 e^{P_j}=e^{X+c_j}e^{\zeta_j}
             \prod_{l<j}e^{-d_l\zeta_l},\qquad
 c_j=\theta_j-\sum_{l<j}d_l\theta_l
     =B_0-\sum_{l\leq j}\log a_l.
\]
The terminal branch has $c_m=-\sum_{l<m}d_l\theta_l$ and the same
last expression. Summing genuine stopped capacities and using the
backward estimates proves for a surviving set $\mathcal S\subseteq
\bigcap_jE_j$ that
\[
 I(\mathcal S)\leq C_H
 \min\{1,e^X\prod_ja_j^{r_j}\},\qquad
 (\prod_ja_j^{r_j})^{\sigma_1}
 =\prod_{j<m}M_j^{\sigma_j}\prod_jp_j^{\sigma_j}.
\]
If a $p_j$ vanishes, the conditional bound makes the finite surviving
support empty, and no logarithm is required. For $m=1$ the direct
full-word capacity proves the assertion. This calculation preserves
the actual descendant lists and is the simple-owner tower used in
the causal proof; hard owners supply their own three conditional
kernels in Lemma~\ref{upper:hard-kernels}.

\subsection{The ordinary rectangle and the same-count root integral}
\label{updetail:ordinary-root}
We prove Proposition~\ref{upper:root-contract}. Write the unread original
suffix as blocks with $h_1>\cdots>h_m\geq2$, $w_i=\log h_i$,
$c_i=h_i-1$, lengths $L_i\geq0$, and $U=\sum_iL_i$.
Let $0<t\leq s<1$, $r=t/s$, in the fixed compact ranges.
With the entire prefix fixed, take arbitrary $D>0$ and $\kappa\in\R$,
and set
\[
 Z=\sum_i(w_iN_i-c_iL_i),\qquad X=\kappa+rZ.
\]
After the positive enlargement, the $N_i$ are independent Poisson
counts with means comparable to $L_i$. The natural means are
$h_i^tL_i$. Conditional on each count there are its original uniform
positions, with all $n_i+1$ gaps. Completed harmonic-cell budgets are
used consistently in $Z$, and their bounded difference from physical
budgets is included in $\kappa$, keeping $X$ unchanged.

Choose a largest original block $M=L_{i_*}$ and put
\[
 A=\sum_{i<i_*}L_i,\quad B=\sum_{i>i_*}L_i,\quad
 B_{\rm act}=\sum_{i>i_*}(w_iN_i-c_iL_i),\quad
 A_{\rm act}=-\sum_{i<i_*}(w_iN_i-c_iL_i).
\]
Then $Z=B_{\rm act}+Z_*-A_{\rm act}$ and $M\leq U\leq mM$.
The two outside sums are independent and retain all their counts.

Three actual rectangles are used before enlargement. A threshold at
the beginning of the suffix assigns its primes freely and bounds
earlier contributions by their full original rectangles, giving
$I\leq Ce^Z$. A threshold within the largest block, with its $j$
smallest prime logarithms retained, gives
$I\leq Ce^{B_{\rm act}}W_{*,N_*}$, where $W_{*,n}$ is the capped
ordered-sum rectangle of \eqref{ordsum:rectangle}. To verify this,
the $h_*-1$ unresolved widths are bounded by the prefix scale times
$1+\sum_{l\leq j}e^{(w_*/c_*)\operatorname{cell}_l}$, whereas
the free later labels contribute $h_*^{-j}$ after normalization.
Minimize over $j$, and use the right endpoint rectangle for the cap
one. The remaining rectangle is $I\leq C$.
Thus, together with a prefix cap $I\leq CDe^X$, the same tuple obeys
\begin{equation}
 I\leq C\widehat V,\qquad
 \widehat V=\min\{1,e^Z,e^{B_{\rm act}}W_{*,N_*},De^X\}.
 \label{updetail:ordinary-envelope}
\end{equation}
For every $n\geq0$ and $\ell\geq0$,
Lemma~\ref{ordsum:rectangle-tail} gives
\[
 \mathbb P(W_{*,n}>e^{-\ell}\mid n)
 \leq\min\left\{1,
 \frac{C(1+\ell)(C_1+Z_*+\ell)_+}{1+n}\right\}.
\]
For negative $\ell$ the event is empty.

First fix all outside counts and restrict $Z$ to an interval
$[E,E+W]$ of fixed width. There are at most $1+W/\log2$ possible
$N_*$. The two full Poisson mode bounds
\begin{equation}
 \sup_n p_\mu(n)\leq C(1+\mu)^{-1/2},\qquad
 \sup_n\frac{p_\mu(n)}{1+n}\leq C(1+\mu)^{-3/2}
 \label{updetail:poisson-modes}
\end{equation}
follow from Stirling near the mode and the elementary Poisson tails;
they include $n=0$ and all noncentral counts. The event
$e^{B_{\rm act}}W_{*,N_*}>e^{-\ell}$ has left reserve
$1+\ell+B_{\rm act}$ and right reserve
$C_1+Z+\ell+A_{\rm act}$. The endpoint rectangle also requires
$Z+\ell\geq0$. Apply the two bounds in
\eqref{updetail:poisson-modes} separately, then average the independent
sides. Since $\mathbb E(B_{\rm act})_+\leq CB$ and
$\mathbb E(A_{\rm act})_+\leq CA$, this gives
\begin{align*}
 &\mathbb P\{Z\in[E,E+W],\ e^Z>e^{-\ell},\
                 e^{B_{\rm act}}W_{*,N_*}>e^{-\ell}\}\\
 &\qquad\leq C(1+U)^{-1/2}
 \min\left\{1,\frac{(1+\ell+B)(1+(E+\ell)_++A)}{1+M}\right\}.
\end{align*}
The root atom and ballot denominator have already been obtained
from the same $N_*$.

Put $H_0=(-\log D-\kappa)/r$ and $w=Z-H_0$. Then
$e^{-sX}=D^se^{-tw}$ and $De^X=e^{rw}$. Positive layer cake in
\eqref{updetail:ordinary-envelope} gives
\begin{align*}
 D^{-s}\mathbb E[e^{-sX}\widehat V]
 =\int_0^\infty e^{-\ell}\,
 \mathbb E\big[e^{-tw}1_{\{w>-\ell/r\}}
 1_{\{Z+\ell\geq0\}}
 1_{\{W_*>e^{-\ell-B_{\rm act}}\}}\big]\,d\ell.
\end{align*}
At fixed outside counts, $a=w+\ell/r>0$ runs along a translated
lattice with spacing $w_*$. Moreover
$Z+\ell=H_0+a-(1/r-1)\ell$, so its nonnegativity forces
$a\geq-H_0$ up to a fixed geometric translation. Uniformly over
the lattice translate and $b\geq0$,
\begin{align*}
 \sum_{a>0,\ a\geq-H_0-C}e^{-ta}&\leq Ce^{-t(H_0)_-},\\
 \sum_{a>0,\ a\geq-H_0-C}e^{-ta}(C+H_0+a+b)_+
 &\leq Ce^{-t(H_0)_-}(1+(H_0)_++b).
\end{align*}
For $H_0<0$, shift by $(H_0)_-$; for $H_0\geq0$, start at zero.
This explains why no polynomial loss in the negative deficit occurs.
Using first the first mode bound and then the second with the ballot
tail bounds the conditional integrand before $e^{-\ell}$ by the
minimum of
\[
 Ce^{s\ell-t(H_0)_-}(1+M)^{-1/2},\qquad
 Ce^{s\ell-t(H_0)_-}(1+M)^{-3/2}
 [1+\ell+(B_{\rm act})_+][1+(H_0)_++(A_{\rm act})_+].
\]
Average the independent sides separately in these two bounds, then
take their minimum. Since $1-s$ has a fixed positive lower bound,
integration against $e^{-(1-s)\ell}$ yields
\begin{equation}
 \mathbb E[e^{-sX}\widehat V]
 \leq CD^se^{-t(H_0)_-}(1+U)^{-1/2}
 \min\left\{1,\frac{(1+B)(1+(H_0)_++A)}{1+M}\right\}.
 \label{updetail:joint-root}
\end{equation}
Missing sides contribute exactly one; an empty suffix is treated by
the deterministic minimum with the same $1+$length conventions.

The outgoing profile needed in the causal theorem is
$J_\beta(X+\log D)$ instead of $De^X$, where
$J_\beta(v)=\min\{1,C_\beta e^v(1+(-v)_+)^\beta\}$ and
$\beta\geq0$ is fixed. If $J_\beta(v)>e^{-\ell}$, then
\[
 v>-L_\beta(\ell),\qquad
 L_\beta(\ell)=\ell+2\beta\log(2+\ell)+C_\beta'.
\]
Indeed $(1+u)^\beta\leq Ce^{u/2}$ first gives
$u=(-v)_+\leq2\ell+C$, which substituted back gives the stated
bound. Now set $q=r(Z-H_0)+L_\beta(\ell)>0$. The tilt becomes
$D^se^{sL_\beta(\ell)}e^{-sq}$ and
$Z+\ell\leq H_0+q/r$. The endpoint cap still forces
$q\geq r(H_0)_-$ up to a fixed translation. The same lattice sums
give $e^{-t(H_0)_-}$ and the same right reserve. The only change in
the layer integration is the weight
\[
 e^{-\ell+sL_\beta(\ell)}
 \leq Ce^{-(1-s)\ell}(2+\ell)^{2s\beta},
\]
whose integral, also multiplied by $1+\ell$, is uniformly finite.
This proves \eqref{updetail:joint-root} with the polynomial profile
and hence Proposition~\ref{upper:root-contract}.

For the natural suffix put
$D_{\rm mean}=\sum_iL_i(h_i^t\log h_i-(h_i-1))\geq0$ and
\[
 \mathcal B_I=(1+U)^{-1/2}
 \min\{1,(1+B)(1+D_{\rm mean}+A)/(1+M)\}.
\]
The preceding bound also gives
\[
 \mathbb E[e^{-sX}\widehat V]
 \leq CD^s\mathcal B_Ie^{-t(H_0)_-}
 \left[1+\frac{(H_0-D_{\rm mean})_+}{1+D_{\rm mean}+A}\right].
\]
The Euler root identity is
$\kappa=-rD_{\rm mean}+\Delta_{\rm pre}$, so
$H_0-D_{\rm mean}=-(\Delta_{\rm pre}+\log D)/r$.
This is the explicit prefix deviation weight used in the final upper
expectation. It remains together with $D^s$ under that expectation;
the conditional estimate allows arbitrary correlation among prefix
costs, counts and geometric sums.
